% https://www.overleaf.com/read/jydxqkkkskzp
% https://github.com/MCG-NKU/NSFC-LaTex
% by Ming-Ming Cheng https://mmcheng.net

%\documentclass[12pt]{article}
\documentclass[12pt,AutoFakeBold]{article}
\usepackage[UTF8]{ctex}
\usepackage{nsfc}
\usepackage{enumitem}

\usepackage{amssymb}

\newcommand{\cmm}[1]{\textcolor[rgb]{0,0.6,0}{CMM: #1}}
\newcommand{\todo}[1]{{\textcolor{red}{\bf [#1]}}}
\newcommand{\myPara}[1]{\paragraph{#1：}}

% yaxing: begin
\usepackage[resetlabels]{multibib}
\newcites{sec}{Reference}
\usepackage{color}

\newcommand{\upcite}[1]{{\textsuperscript{\cite{#1}}}}
\newcommand{\sprod}[2]{\left\langle #1,#2 \right\rangle}
\newcommand{\minisection}[1]{\vspace{0.04in}\noindent {\bf #1}\ \ }


%  \iffalse
\newcommand{\revision}[1]{{\color{blue} #1}}
\graphicspath{{figures/}}


\begin{document}

\begin{figure}[h]
    \centering
    \includegraphics[width=.7\linewidth]{figures/logo.jpeg}
\end{figure}

\vspace{1cm}
\hrulefill
    \begin{center}
        \Huge\textbf{《人工智能导论》探究报告\\ 图像超分\textcolor{red}{为例} }
     
    \end{center}   

\hrulefill
\vspace{1cm}

\begin{table}[h]
    \centering
    \large
    \begin{tabular}{ll}
    \textbf{姓名：XXX   \quad  学号:0000000} \quad 论文1  \\
    \textbf{姓名：XXX    \quad  学号:0000000} \quad 论文2 \\
    \textbf{姓名：XXX    \quad  学号:0000000} \quad 论文3 \\
    \end{tabular}
\end{table}

\newpage

%%%%%%%%% TITLE

\title{\kaishu 作业正文}

\maketitle




\section{问题描述（\revision{文字描述与形式化定义，组员共同完成}）}


\section{核心内容：}


% \section{核心内容（\revision{三篇论文三位组员分别负责，拓展阅读论文组员共同选题和完成，描述论文动机，方法，实验结果等关键内容}）：}

\subsection{\textbf{Real-ESRGAN: Training Real-World Blind Super-Resolution with Pure Synthetic Data}\cite{wang2021realesrgantrainingrealworldblind}}

\subsubsection{论文动机}

% \paragraph{123}

在图像超分辨率（Super-Resolution, SR）领域，深度学习方法已经取得了显著的进展。早期如 SRGAN、ESRGAN 等方法主要依赖于人工合成的低分辨率-高分辨率（LR-HR）图像对进行训练，这些图像对通常通过对高清图像进行 Bicubic（双三次） 等方式降采样获得。

然而，在现实中，低分辨率图像往往受到多种复杂因素影响（如模糊、噪声、压缩伪影、色彩失真等），甚至是多种复杂退化过程的叠加，比如例如相机的成像系统、图像编辑和互联⽹传输等等，而不是仅仅的简单下采样缩放。
因此，使用合成的 LR-HR 图像对进行训练，无法全面反映真实世界中的图像退化特性。这种合成退化与实际应用中的低分辨率图像存在较大差距，极大地限制了模型的泛化能力和实际应用效果。

可是，真实数据集的缺失无可避免。一方面，收集真实世界的低分辨率图像和对应的高分辨率图像对非常困难且成本高昂，在之前RealSR\cite{RealSR}的工作中，虽然提出了一种通过用同一相机不同焦距下拍摄来构建真实世界图像超分辨率数据集的方法，但数据集制造成本仍然昂贵，并且规模和多样性有限，难以应对多种实际场景。另一方面，获取同一场景、同一内容下高分辨率（HR）和低分辨率（LR）图像对，技术和成本门槛高，且存在对齐误差、光照变化等问题。

为了应对真实数据集缺失、合成数据集泛化能力差的挑战，本文提出了一种更加复杂、更加多样化的退化建模方法，来模拟真实场景下的多元退化，提升训练数据的多样性和代表性，缓解真实 LR-HR 对数据稀缺的问题。同时退化的过程更加贴近现实，模拟现实中会出现的振铃和伪影，使训练出的模型能更好地泛化到真实场景。

此外，先前ESRGAN\cite{ESRGAN}中的判别器结构和感知损失设计容易在真实退化场景下学到伪特征，导致重建结果过度锐化或出现假纹理，而感知损失与对抗损失的组合又容易引发梯度爆炸或者消失，使得训练不稳定，影响模型落地。

所以本文的工作还把判别器中的VGG结构换成了U-Net结构，同时加入了谱归一化的操作，让训练过程更加稳定。

总的来说，本文旨在通过设计更加贴近真实世界的退化建模方法，构建兼具代表性和多样性的训练数据集，并突破以往方法的局限性。论文尝试仅依赖纯合成数据进行模型训练，但通过更合理的退化流程和网络设计，使得训练出的模型能够在真实场景中依然具备优异的泛化能力和实际应用效果。


总的来说，本文旨在通过设计更加贴近真实世界的退化建模方法，来构建更加贴近现实、兼具代表性和多样性的数据集。同时，论文尝试仅依赖纯合成数据进行模型训练，但通过更合理的退化流程和网络设计，使得训练出的模型能够在真实场景中依然具备优异的泛化能力和实际应用效果。

% 如何使用插入图像：


% \begin{figure*}[h]
%     \centering
%     \includegraphics[width=1\linewidth]{figures/teaser.jpg}  
%     %\vspace{-6mm}
%     \caption{\small 示意图}
%     \label{fig:framework}
%     %\vspace{-15pt}
% \end{figure*}



\subsubsection{方法-退化模型部分}

本论文工作的方法主要分为两部分：退化模型和Real-ESRGAN模型。我们先来介绍一些前置知识以及本文提出的高阶退化模型：

%----------------------------------
\paragraph{经典退化模型}

经典的退化模型主要是对高分辨率图像$y$先用一个模糊核$k$进行卷积操作，然后执行下采样，加上噪声$n$之后，最后再经过一个JPEG压缩，得到低分辨率图像。
\[
\mathbf{x} = \mathcal{D}(y) = \left[ ( y \circledast k ) \downarrow_r + n \right]_{JPEG}
\]

{\small 这里的$\mathcal{D}(y)$表示退化过程；$\downarrow_r$代表比例因子为r的下采样操作；$( y \circledast k )$表示对图像$y$,用模糊核$k$进行卷积操作。}


%----------------------------------
\paragraph{模糊（Blur）}
在经典的退化模型中，通常将模糊退化过程用原图像与某个线性模糊核进行卷积的方式实现。模糊核经常选择各向异性或者各向同性的高斯模糊核。对于尺寸为 \(2t + 1\) 的高斯模糊核 \(k\)，其在 \((i, j) \in [-t, t]\) 范围内的元素为从高斯分布采样，形式如下：

\[
k(i,j) = \frac{1}{N} \exp \left( -\frac{1}{2} C^{T} \Sigma^{-1} C \right),
\quad C = [i, j]^{T}
\]

{\small 其中，\(\Sigma\) 为协方差矩阵；\(C\) 为空间坐标向量；\(N\) 为归一化常数。}

\begin{figure}[h]
    \centering
    \includegraphics[width=0.8\linewidth]{截屏2025-05-29 00.04.18.png}
    \caption{几种高斯核模糊图像示例}
    \label{fig:2}
\end{figure}

% 协方差矩阵还可以进一步表示为：
% \[
% \Sigma = R \begin{bmatrix}
% \sigma_1^2 & 0 \\
% 0 & \sigma_2^2
% \end{bmatrix}
% R^{T}, \quad 
% \]
% \[
% = \left[\begin{array}{cc}
% \cos \theta & -\sin \theta \\
% \sin \theta & \cos \theta
% \end{array}\right]
% \begin{bmatrix}
% \sigma_1^2 & 0 \\
% 0 & \sigma_2^2
% \end{bmatrix}
% \left[\begin{array}{cc}
% \cos \theta & \sin \theta \\
% -\sin \theta & \cos \theta
% \end{array}\right], \quad 
% \]
% 其中，\(\sigma_1\) 和 \(\sigma_2\) 分别是沿两个主轴的标准差（即协方差矩阵的特征值）；\(\theta\) 为旋转角度。当 \(\sigma_1 = \sigma_2\) 时，\(k\) 为各向同性高斯模糊核，否则为各向异性模糊核。


即使高斯模糊核被广泛使用，但是它可能无法很好的近似真实的相机模糊，为了让模糊核更加多样化，本文的退化模型选择使用广义高斯模糊核和台形分布（中间部分权重相等，外部权重逐渐减小）。其概率密度函数形式如下：

\[
f(x) = \frac{1}{N} \exp \left( -\frac{1}{2} (C^{T} \Sigma^{-1} C)^{\beta} \right) \text{ or }
f(x) = \frac{1}{N} \frac{1}{1 + (C^{T} \Sigma^{-1} C)^{\beta},} \quad
\]
{\small 其中，\(\beta\) 为形状参数。}


%----------------------------------
\paragraph{添加噪声（Add Noise）}
在 Real-ESRGAN 的退化过程中，噪声是一类重要且常见的图像退化因素。文中主要考虑了两种类型的噪声：加性高斯噪声（Additive Gaussian Noise）和泊松噪声（Poisson Noise）。以更全面地模拟真实拍摄和传输过程中产生的不同噪声类型。

\subparagraph{加性高斯噪声}具有与高斯分布相同的概率密度函数。噪声强度由高斯分布的标准差（即 $\sigma$ 值）控制。当对RGB图像的每个通道分别独立采样噪声时，合成的噪声属于彩色噪声（color noise）；而如果对三个通道应用相同的噪声样本，则为灰度噪声（gray noise）\cite{RealESRGAN, BM3D}。高斯噪声可以有效模拟传感器热噪声、电子干扰等随机噪声现象，增强模型对噪声的鲁棒性。

\subparagraph{泊松噪声}遵循泊松分布。它通常用于近似模拟由统计量子涨落引起的传感器噪声——即在给定曝光水平下，光子数的变化所导致的噪声。泊松噪声的强度与图像像素本身的强度成正比，并且不同像素位置的噪声相互独立。这种噪声类型常见于低光照、医学成像等实际场景，使退化模型更加贴近真实。


\paragraph{下采样（Downsampling）}
在超分辨率（SR）任务中，合成低分辨率（LR）图像的一个基本步骤就是下采样操作。更广义地说，调整图像大小不仅包括下采样，也包括上采样。图像大小的调整方式有多种常见算法，例如最近邻插值、区域调整（Area）、双线性插值（Bilinear）和双三次插值（Bicubic）。不同的调整算法会对图像产生不同的影响：部分方法（如双线性、双三次）通常会带来一定程度的模糊效果，而有些方法（如最近邻）可能会引入过冲伪影或导致图像过锐。

为了使合成的退化图像更加多样和复杂，Real-ESRGAN 在退化过程中采用了“随机调整大小”策略，即从上述多种调整大小算法中随机选择一种进行操作。但具体而言，考虑到最近邻插值容易导致像素错位等问题，Real-ESRGAN 排除了最近邻插值，仅在区域调整、双线性插值和双三次插值三种算法中进行随机选择。这样设计能够模拟现实世界中因不同成像设备、压缩算法或编辑软件带来的多样化降质效应，从而提升模型在真实场景下的泛化能力和鲁棒性。



\paragraph{JPEG 压缩（JPEG Compression）}


[43]。


JPEG 压缩是一种广泛应用的有损数字图像压缩技术。在 JPEG 压缩流程中，首先将图像从 RGB 空间转换为 YCbCr 颜色空间，并对色度通道（Cb、Cr）进行下采样，以减小数据量。随后，图像会被分割为 $8 \times 8$ 的小块，每个块都通过二维离散余弦变换（DCT）将像素值转换到频域。DCT 系数经过量化操作以进一步压缩信息。

JPEG 压缩往往会引入明显的块状伪影（blocking artifacts），影响图像的主观质量。压缩图像的最终质量由质量因子 $q \in [0, 100]$ 控制，$q$ 值越低，压缩率越高但图像质量越差；$q$ 值越高，保留更多细节但文件体积较大。

在 Real-ESRGAN 的退化流程中，采用了 PyTorch 实现的 DiffJPEG 算法[32]，以在训练过程中灵活地对图像进行可微分的 JPEG 压缩操作，从而更真实地模拟实际使用场景下的压缩噪声和伪影。


\paragraph{振铃和过冲伪影}
\textbf{振铃伪影}常见于图像中锐利过渡（如边缘）附近，表现为虚假的边界、条带或“幽灵”效应。\textbf{过冲伪影}通常与振铃伪影同时出现，表现为在边缘跳变处像素值异常跃升，使得边缘部分产生过度增强的效果。  

这些伪影的主要成因是信号被带限，高频成分被截断，导致复原或处理后无法真实还原原始细节。这类伪影在图像锐化、JPEG压缩等处理中非常常见。实际图像中常见于过度锐化和压缩后。

如图\ref{fig:3}所示，一些真实图像样本中出现了明显的振铃和过冲伪影；它们在边缘附近表现为明显的条带和亮/暗带。

\begin{figure}[h]
    \centering
    \includegraphics[width=0.8\linewidth]{截屏2025-05-29 16.26.52.png}
    \caption{振铃和过冲伪影示例}
    \label{fig:3}
\end{figure}


为模拟这类伪影，Real-ESRGAN 在训练数据的合成过程中引入了Sinc滤波器。Sinc滤波器是一种理想的低通滤波器，能够截断高于某一截止频率（ωc）的所有高频分量，从而合成出类似振铃和过冲的伪影。

sinc滤波器核的表达式如下：  
\[
k(i,j) = \frac{\omega_c}{2\pi \sqrt{i^2 + j^2}} J_1 \left( \omega_c \sqrt{i^2 + j^2} \right),
\]

其中 $(i, j)$ 为滤波器核的坐标，$\omega_c$ 为截止频率，$J_1$ 为一阶第一类贝塞尔函数。  

在 Real-ESRGAN 的退化流程中，sinc滤波器被用于模糊过程和作为最后一步的滤波操作。为了进一步增加退化的多样性，最后一步的 sinc 滤波与 JPEG 压缩的顺序是随机的——即有些图像会先锐化（产生过冲/振铃）再压缩，有些则先压缩再锐化。这样既可以模拟先过度锐化后压缩的伪影，也能模拟先压缩后锐化的情形，极大丰富了退化空间。

\paragraph{高阶退化模型}

由于用上面提到的经典退化模型难以模拟现实世界中复杂的图像退化过程，Real-ESRGAN 进一步提出了高阶退化模型。该模型通过将上述经典退化模型中的每个步骤进行多次迭代，来模拟更复杂的退化过程。

一个n阶的退化模型会重复n次经典退化模型的过程，但是每次的超参数可能会不同。具体而言，n阶退化模型的形式如下：
\[
x = \mathcal{D}^n(y) = (\mathcal{D}_n \circ \cdots \circ \mathcal{D}_2 \circ \mathcal{D}_1)(y).
\]

{\small 其中，\(\mathcal{D}^{n}(y)\) 表示n阶退化模型对图像\(y\)的退化过程。}

根据经验，Real-ESRGAN采⽤⼆阶退化过
程，可以在保持简单性的同时解决⼤多数实际问题。图\ref{fig:1}描绘了纯合成数据⽣成流程的整体流程。

\begin{figure}[h]
    \centering
    \includegraphics[width=1\linewidth]{截屏2025-05-28 22.34.31.png}
    \caption{退化模型流程图}
    \label{fig:1}
\end{figure}

\subsubsection{方法-训练模型部分}

Real-ESRGAN是基于前人的工作上面改进的，在介绍 Real-ESRGAN 所做的改进之前，有必要简要回顾一下前身——SRGAN模型\cite{SRGAN}和ESRGAN模型\cite{ESRGAN}。

\paragraph{\textbf{SRGAN}}
是一个开创性的图像超分辨率生成对抗网络，SRGAN 的主要创新在于引入了生成对抗网络（GAN）结构，将超分辨率重建任务转化为一个生成器与判别器相互博弈的过程，从而获得更高感知质量的高分辨率图像。

SRGAN 的模型结构主要包括以下两个部分,结构如图\ref{fig:4}所示：

\begin{enumerate}
    \item 生成器（Generator）
    
    生成器负责将低分辨率（LR）图像转换为高分辨率（HR）图像。其网络结构以深度残差网络为基础，并引入了跳跃连接，以增强网络的表达能力和训练的稳定性。
    
    \item 判别器（Discriminator）

    判别器的任务是判断输入图像是真实的高分辨率图像，还是由生成器合成出来的超分辨率图像。通过与生成器的对抗训练，判别器不断提升区分能力，促使生成器生成更加逼真的细节。
\end{enumerate}

\begin{figure}[h]
    \centering
    \includegraphics[width=1\linewidth]{截屏2025-05-29 16.42.24.png}
    \caption{SRGAN模型}
    \label{fig:4}
\end{figure}

    然后，SRGAN 的损失函数由三部分加权组成：
    \begin{enumerate}
        \item \textbf{内容损失}

        通常采用感知损失，即在 VGG 网络的特征空间中计算 LR 和 HR 图像的差异，从而提升重建图像的主观感知质量。
        
        \item \textbf{对抗损失}

        来自判别器的反馈，鼓励生成器生成更具真实感的纹理和细节。这里SRGAN的反馈是判别器返回的一个输入为真高分辨率图的概率。
        
        \item \textbf{像素损失}

        即传统的像素级均方误差（MSE）损失，用以保证基础结构的准确还原。
        
    \end{enumerate}


\paragraph{\textbf{ESRGAN}}是在 SRGAN 基础上提出的改进方法，主要针对 SRGAN 在真实世界图像超分辨率任务中的不足之处进行了优化。ESRGAN 的主要改进包括：



\begin{figure}[h]
    \centering
    \includegraphics[width=1\linewidth]{截屏2025-05-29 17.05.58.png}
    \caption{ESRGAN网络改进}
    \label{fig:enter-label}
\end{figure}

\begin{figure}[h]
    \centering
    \includegraphics[width=1\linewidth]{截屏2025-05-29 17.04.04.png}
    \caption{相对判别器改进}
    \label{fig:5}
\end{figure}

\subsubsection{实验结果}

% 如何插入表格：

% \begin{table}[h]
%   \centering
%   \begin{tabular}{@{}lc@{}}
%     \hline    \hline

%     Method & Frobnability \\
%     \hline
%     Theirs & Frumpy \\
%         \hline

%     Yours & Frobbly \\
%         \hline

%     Ours & Makes one's heart Frob\\
%     \hline    \hline

%   \end{tabular}
%   \caption{Results.   Ours is better.}
%   \label{tab:example}
% \end{table}
\subsection{规定论文题目2}

\subsection{规定论文题目3}

\subsection{自选论文题目 （\revision{组员共同完成}）}




\section{思考与理解（\revision{联系与区别、尚未解决的问题、未来研究趋势等，组员共同完成}）：}

\subsection{联系与区别}

\subsection{尚未解决的问题}


\subsection{未来研究趋势}

% 如何添加参考文献：

% 图像到图像转换基本思想是将给定源域图像映射到目标域，在此过程中输出图像具有两个特点：(1)保持输入图像结构、姿势信息；(2)获取目标域图像风格、属性、外表特征。图像到图像转换最早追溯到图像类比\upcite{hertzmann2001image}开始， Hertzmann等人\upcite{efros1999texture}提出单个输入-输出训练图像对的非参数纹理模型。 最近的方法\upcite{long2015fully}使用输入输出示例的数据集来学习使用卷积神经网络(CNN)的参数翻译函数。针对输出为RGB的图像到图像转换，近年来相关研究主要集中于基于生成对抗网络(GAN)的学习方法，其从数据要求角度通常分为二类：监督方式图像到图像转换和非监督方式图像到图像转换。监督方式图像到图像转换中，Phillip等人\upcite{pix2pix2017}基于U-Net构建一种包含编码器-解码器的生成器架构：Pix2pix，中科院团队\upcite{song2018geometry}在图像转换中引入关键点信息


\newpage
{\small
\bibliography{longstrings,Cmm}
\bibliographystyle{unsrt}
}

\end{document}
